\chapter{Guida Utente}

\section{Scopo del gioco}

Il giocatore dovrà risolvere una serie di puzzle tramite la stesura di semplici
codici Assembler-like componibili trascinando le istruzioni che si vuole
utilizzare nella zona a destra della lista delle istruzioni fornite dal
livello. Quando verrà avviata l'esecuzione del codice scritto, la navicella
inizierà a svolgere il programma eseguendo una ad una e in ordine le istruzioni
dell'utente.  Per ogni livello l'utente potrà provare a scrivere il codice che
fa uso del minor numero di istruzioni possibile o che risolve il problema col
minor numero di passi possibile, nel caso uno degli obbiettivi venga raggiunto
l'utente ne verrà notificato dall'applicazione.

\section{Interfaccia grafica}

\begin{figure}[H]
\centering
\includegraphics[width=0.8\textwidth]{images/interface}
\caption{Schermata di gioco di un livello}
\end{figure}

\begin{enumerate}
\def\labelenumi{\arabic{enumi}.}
\itemsep0em
\item
  Torna al menù della scelta dei livelli
\item
  Mostra l'obbiettivo del livello
\item
  Riduci la velocità di esecuzione del codice
\item
  Aumenta la velocità di esecuzione del codice
\item
  Cancella tutte le istruzioni inserite
\item
  Copia il codice scritto fino ad ora
\item
  Incolla un codice copiato
\item
  Annulla l'ultima modifica al codice
\item
  Ripeti l'ultima modifica annullata al codice
\item
  Avvia l'esecuzione ininterrotta del codice scritto fino ad ora
\item
  Avvia l'esecuzione step-by-step del codice
\item
  Nastro contenente l'input del programma
\item
  Area di memoria del programma composta da singole celle identificate
  da un indice
\item
  Navicella che eseguirà le istruzioni
\item
  Nastro contenente l'output del programma
\item
  Istruzioni utilizzabili dall'utente
\item
  Programma composto dall'utente
\end{enumerate}

\section{Istruzioni disponibili}

\begin{figure}[H]
    \centering
    \begin{minipage}[t]{0.49\textwidth}
        \includegraphics[width=\textwidth]{images/instructions/input}
        \caption{
            La navicella prende un valore dal nastro di input (scartando eventuali
            valori precedentemente tenuti)
        }
    \end{minipage}\hfill
    \begin{minipage}[t]{0.49\textwidth}
        \includegraphics[width=\textwidth]{images/instructions/output}
        \caption{
            La navicella inserisce il valore tenuto nel nastro di output
        }
    \end{minipage}
\end{figure}

\begin{figure}[H]
    \centering
    \begin{minipage}[t]{0.32\textwidth}
        \includegraphics[width=\textwidth]{images/instructions/jump}
        \caption{
            Salta al target indicato
        }
    \end{minipage}\hfill
    \begin{minipage}[t]{0.32\textwidth}
        \includegraphics[width=\textwidth]{images/instructions/jumpzero}
        \caption{
            Salta al target indicato solo se il valore tenuto dalla navicella è
            0
        }
    \end{minipage}\hfill
    \begin{minipage}[t]{0.32\textwidth}
        \includegraphics[width=\textwidth]{images/instructions/jumpneg}
        \caption{
            Salta al target indicato solo se il valore tenuto dalla navicella è
            negativo
        }
    \end{minipage}
\end{figure}

\begin{figure}[H]
    \centering
    \begin{minipage}[t]{0.49\textwidth}
        \includegraphics[width=\textwidth]{images/instructions/copyfrom}
        \caption{
            La navicella prende una copia (scartando eventuali valori
            precedentemente tenuti) del valore nella cella di memoria indicata
            dal numero a destra
        }
    \end{minipage}\hfill
    \begin{minipage}[t]{0.49\textwidth}
        \includegraphics[width=\textwidth]{images/instructions/copyto}
        \caption{
            Inserisce nella cella di memoria indicata dal numero a destra una
            copia del valore tenuto dalla navicella
        }
    \end{minipage}
\end{figure}

\begin{figure}[H]
    \centering
    \begin{minipage}[t]{0.49\textwidth}
        \includegraphics[width=\textwidth]{images/instructions/add}
        \caption{
            Somma al valore tenuto dalla navicella il valore che si trova nella
            cella di memoria indicata dal numero a destra
        }
    \end{minipage}\hfill
    \begin{minipage}[t]{0.49\textwidth}
        \includegraphics[width=\textwidth]{images/instructions/sub}
        \caption{
            Sottrae dal valore tenuto dalla navicella il valore che si trova
            nella cella di memoria indicata dal numero a destra
        }
    \end{minipage}
\end{figure}

\begin{figure}[H]
    \centering
    \begin{minipage}[t]{0.49\textwidth}
        \includegraphics[width=\textwidth]{images/instructions/incr}
        \caption{
            Incrementa di una unità il valore numerico che si trova nella cella
            di memoria indicata dal numero a destra
        }
    \end{minipage}\hfill
    \begin{minipage}[t]{0.49\textwidth}
        \includegraphics[width=\textwidth]{images/instructions/decr}
        \caption{
            Decrementa di una unità il valore numerico che si trova nella cella
            di memoria indicata dal numero a destra
        }
    \end{minipage}
\end{figure}
